%!TEX TS-program = xelatex
%!TEX encoding = UTF-8

% Unit pembaca Bahasa Indonesia untuk Methods of Algebra, Volume 1.
% Teks sumber: Copyright 2018 Wen-Wei Li, CC BY 4.0.
% Terjemahan ini dilisensikan menurut CC BY 4.0.
% Provenans produksi: OpenAI Codex gpt-5.6-sol, Ultra.

\documentclass[
	colors = true,
	coverpage = coverpage-id-unit-034.tex,
	fontsetup = font-setup-id.tex,
	titlesetup = titles-setup-id.tex,
	CJKthechapter = false
]{AJbook}

% Target-only Indonesian labels for source-class reader interface elements.
% The upstream AJbook.cls remains byte-identical to the frozen source closure.
\renewtcolorbox{wenxintishi}{
	breakable,
	enhanced,
	width = \textwidth,
	colback = white, colbacktitle = white,
	colframe = gray!50, boxrule=0.2mm,
	coltitle = black,
	fonttitle = \sffamily,
	attach boxed title to top left = {yshift=-\tcboxedtitleheight/2, xshift=\tcboxedtitlewidth/4},
	boxed title style = {boxrule=0pt, colframe=white},
	before skip = 0.5cm,
	top = 3mm,
	bottom = 3mm,
	title = {Petunjuk membaca}
}


\makeatletter
\AtBeginDocument{\gdef\mkheader@lemma{\setparms@lemma\@thm{lemma}{theorem}{Lema}}}
\makeatother

\setCJKmainfont[
	Path=../fonts/,
	BoldFont=NotoSansCJKsc-Medium.otf,
	ItalicFont=NotoSansCJKsc-Regular.otf
]{NotoSansCJKsc-Regular.otf}

\DefineBibliographyStrings{english}{
	in = {dalam},
	editor = {penyunting},
	byeditor = {disunting oleh},
	backrefpage = {dirujuk pada hlm\adddot},
	backrefpages = {dirujuk pada hlm\adddot}
}

\newcommand{\BuildDateID}{27 Agustus 2026}
\newcommand{\BuildDateShortID}{2026-08-27}
\AtBeginDocument{\renewcommand{\chaptermark}[1]{\markboth{Bab \thechapter\quad #1}{}}}
\AtBeginDocument{\let\cleardoublepage\clearpage}

\usepackage{mathrsfs}
\usepackage{stmaryrd}
\SetSymbolFont{stmry}{bold}{U}{stmry}{m}{n}
\usepackage{array}
\usepackage{makecell}
\usepackage{tikz-cd}
\usetikzlibrary{positioning, patterns, calc, matrix, shapes.arrows, shapes.symbols}
\usepackage{braids}
\usepackage{tqft}
\usepackage{ytableau}
\usepackage{pgfplots}
\pgfplotsset{compat=1.18}
\usepackage{myarrows}
\usepackage{mycommand}

\usepackage[noautomatic]{imakeidx}
\makeindex[
	columns=2,
	program=makeindex,
	intoc=true,
	title={Indeks Istilah}]
\makeindex[
	columns=3,
	program=makeindex,
	intoc=true,
	name=sym1,
	title={Indeks Simbol}]

\usepackage{xr}
\usepackage[unicode, bookmarksnumbered]{hyperref}
\externaldocument{unit-034-crossrefs}[]
\let\UnitThirtyFourOriginalRef\ref
\newcommand{\UnitThirtyFourMarkExternalRef}[1]{%
	\expandafter\def\csname unit034external:#1\endcsname{1}}
\newcommand{\UnitThirtyFourRef}[1]{%
	\ifcsname unit034external:#1\endcsname
		\begin{NoHyper}\UnitThirtyFourOriginalRef{#1}\end{NoHyper}%
	\else
		\UnitThirtyFourOriginalRef{#1}%
	\fi}
\UnitThirtyFourMarkExternalRef{sec:free-group}
\UnitThirtyFourMarkExternalRef{eg:complete-cocomplete}
\UnitThirtyFourMarkExternalRef{prop:limit-buildingblocks}
\UnitThirtyFourMarkExternalRef{eg:categories}
\UnitThirtyFourMarkExternalRef{def:filtrant-poset}
\UnitThirtyFourMarkExternalRef{def:filtrant-cat}
\UnitThirtyFourMarkExternalRef{prop:coset-decomp}
\UnitThirtyFourMarkExternalRef{prop:completion-ring-characterization}
\UnitThirtyFourMarkExternalRef{sec:filters}
\UnitThirtyFourMarkExternalRef{sec:ring-limits}
\hypersetup{
	pdfauthor = {Wen-Wei Li},
	pdftitle = {Metode Aljabar, Jilid 1: Arsitektur Dasar - Unit 34: Limit dan Pelengkapan Grup},
	pdfsubject = {Terjemahan Bahasa Indonesia independen; Bagian 4.10 lengkap},
	pdfkeywords = {aljabar, teori grup, limit proyektif, grup topologis, grup profinit, pelengkapan, bilangan bulat p-adik, modul Tate, id-ID},
	pdflang = {id-ID},
	CJKbookmarks = true}

\addbibresource{Al-jabr.bib}

\begin{document}
	\hypersetup{pageanchor=false}
	\pagestyle{empty}
	\input{coverpage-id-unit-034.tex}
	\pagestyle{fancy}
	\pagenumbering{roman}
	\setcounter{page}{1}

	\mainmatter
	\hypersetup{pageanchor=true}
	\raggedbottom
	\newgeometry{
		centering,
		textwidth=142mm,
		textheight=198mm,
		headheight=5ex,
		headsep=5ex}
	\setlength{\headwidth}{\textwidth}
	\setcounter{chapter}{4}
	\setcounter{section}{9}
	% Six numbered displays precede Section 4.10 in the complete chapter.
	\setcounter{equation}{6}
	\begingroup
	\setlength{\emergencystretch}{3em}
	\setstretch{1.16}
	\clubpenalty=10000
	\widowpenalty=10000
	\displaywidowpenalty=10000
	\brokenpenalty=10000
	\let\ref\UnitThirtyFourRef
	% The build script extracts this exact span from the hash-gated canonical target.
	\input{unit-034-canonical-span.tex}
	\endgroup

	\begingroup
	\setlength{\emergencystretch}{3em}
	\par\bigskip
	\phantomsection
	\pdfbookmark[0]{Daftar Pustaka}{unit034-bibliography}
	{\Large\sffamily\bfseries Daftar Pustaka\par}\smallskip
	\begingroup
	\small
	\setlength{\bibitemsep}{0.25\baselineskip}
	\printbibliography[heading=none]
	\endgroup
	\let\clearpage\relax
	\let\cleardoublepage\relax
	\IfFileExists{\jobname.ind}{
		\begingroup
		\renewenvironment{theindex}
			{\par\smallskip
			 \phantomsection
			 \pdfbookmark[0]{Indeks Istilah}{unit034-term-index}
			 \noindent{\small\sffamily\bfseries Indeks Istilah.}\enspace
			 \small
			 \let\UnitThirtyFourIndexPar\par
			 \let\par\relax
			 \def\item{\unskip\hspace{1.25em}\penalty0}
			 \let\indexspace\relax}
			{\let\par\UnitThirtyFourIndexPar\par}
		\input{\jobname.ind}
		\endgroup
	}{}
	\IfFileExists{sym1.ind}{
		\begingroup
		\renewenvironment{theindex}
			{\par\smallskip
			 \phantomsection
			 \pdfbookmark[0]{Indeks Simbol}{unit034-symbol-index}
			 \noindent{\small\sffamily\bfseries Indeks Simbol.}\enspace
			 \small
			 \let\UnitThirtyFourSymbolIndexPar\par
			 \let\par\relax
			 \def\item{\unskip\hspace{1.25em}\penalty0}
			 \let\indexspace\relax}
			{\let\par\UnitThirtyFourSymbolIndexPar\par}
		\input{sym1.ind}
		\endgroup
	}{}
	\endgroup
\end{document}
